\documentclass[11pt]{article}

\usepackage[margin=1in]{geometry}
\usepackage{booktabs}

\title{EE 5103 / EE 4953 --- Final Programming Project\\
\large Circuit Analysis Library in C++}
\author{Jordan Cavlovic}
\date{May 2026}

\begin{document}

\maketitle

\section{Problem Description and Motivation}

The goal of this project is to build a small C++ library that represents
simple DC circuits, computes their equivalent resistance, the current
drawn from the source, the voltage drop across each component, and the
power dissipated. The user can construct resistors, capacitors, inductors,
and ideal voltage sources, group them into series or parallel circuits,
combine entire circuits with the \texttt{+} and \texttt{*} operators, and
save or load circuits as text files. At DC steady state an inductor is
treated as a perfect wire and a capacitor as an open circuit, so the
analysis reduces to pure resistance, Ohm's law ($V = IR$), and the power
law ($P = I^{2}R$).

Hand-computing voltage drops and dissipated power becomes tedious past a
few elements, and a small object-oriented library is a much better fit:
each element type knows how to compute its own resistance, voltage drop,
and power, and a container class composes them. The project also exercises
a broad set of modern C++ features in a single, coherent codebase.

\section{Example Run}

The driver prompts the user for a source voltage $V_s$ and three resistor
values, builds a series circuit, and reports the total resistance, the
current, and the voltage drop and power for each element. The same source
voltage is then applied to a parallel network and to series/parallel
compositions built with the \texttt{+} and \texttt{*} operators.

\begin{table}[h]
\centering
\caption{Series-circuit summary for three sample runs ($R_{tot}$ is the
sum of the three resistors, $I = V_s / R_{tot}$, $P_{tot} = V_s^{2}/R_{tot}$).}
\begin{tabular}{@{}rrrrrrr@{}}
\toprule
$V_s$ (V) & $R_1$ ($\Omega$) & $R_2$ ($\Omega$) & $R_3$ ($\Omega$)
& $R_{tot}$ ($\Omega$) & $I$ (A) & $P_{tot}$ (W) \\
\midrule
12 & 100 & 220 & 470 & 790 & 0.01519 & 0.182 \\
24 & 100 & 200 & 300 & 600 & 0.04000 & 0.960 \\
 5 &  10 &  20 &  30 &  60 & 0.08333 & 0.417 \\
\bottomrule
\end{tabular}
\end{table}

\begin{table}[h]
\centering
\caption{Per-element voltage drop and dissipated power for the first run
($V_s = 12$~V, series).}
\begin{tabular}{@{}lrrr@{}}
\toprule
Element & $R$ ($\Omega$) & $V_{drop}$ (V) & $P$ (W) \\
\midrule
$R_1$ & 100 & 1.519 & 0.023 \\
$R_2$ & 220 & 3.342 & 0.051 \\
$R_3$ & 470 & 7.139 & 0.108 \\
\midrule
Total & 790 & 12.000 & 0.182 \\
\bottomrule
\end{tabular}
\end{table}

\section{Software Design and Class Structure}

The design follows the \emph{composite pattern}: a circuit is itself a
component, so a single recursive type system represents both leaf elements
and composite ones. \texttt{Component} is an abstract base class that
declares the pure virtual method \texttt{resistance()} together with
\texttt{type()}, \texttt{clone()}, and \texttt{print()}, and stores the
shared label and value. It also provides virtual default implementations
of \texttt{voltageDrop(I)} ($V = IR$) and \texttt{power(I)} ($P = I^{2}R$),
which subclasses override only when the DC steady-state behavior departs
from the resistance-only formula. The four leaf classes ---
\texttt{Resistor}, \texttt{Capacitor}, \texttt{Inductor}, and
\texttt{VoltageSource} --- each report their DC resistance: $R$ for a
resistor, $0$ for an inductor (a wire), $\infty$ for a capacitor (an open),
and $0$ for an ideal voltage source. Inductor, Capacitor, and
VoltageSource also override \texttt{voltageDrop} and \texttt{power} to
encode the corresponding DC behavior (zero drop and dissipation for the
storage elements, $-V$ and $-V\,I$ for the source so it is reported as
supplying rather than dissipating power).

\texttt{Circuit} also derives from \texttt{Component}, so it can be
embedded inside another circuit transparently. It owns a vector of
children and a series/parallel topology, walks its children to compute the
combined resistance (sum for series, reciprocal-of-reciprocals for
parallel), and exposes an \texttt{analyze(os, $V_s$)} method that prints
the total resistance, the applied voltage and current, the per-component
voltage drop and dissipated power, and the total power. \texttt{analyze}
recurses into nested sub-circuits via \texttt{dynamic\_cast}, applying the
correct branch voltage in each case. \texttt{Circuit} also implements
\texttt{operator+} for series combination, \texttt{operator*} for parallel
combination, and \texttt{save()} / \texttt{load()} for text-file
persistence. \texttt{ComponentFactory} is a small helper class with a
static \texttt{create()} method that maps a type-name string to the
matching concrete object, keeping \texttt{Circuit::load} free of any
knowledge of the leaf classes. The driver in \texttt{main.cpp} wires
everything together. Every node in the tree exposes the same interface,
so client code does not need to distinguish between a single resistor and
a nested sub-circuit.

\section{Discussion of C++ Concepts Used}

\paragraph{Inheritance and runtime polymorphism.}
\texttt{Component} declares pure virtual methods that each concrete subclass
overrides. \texttt{Circuit} stores its children as
\texttt{std::unique\_ptr<Component>}, so a single virtual call dispatches
to the right implementation regardless of whether the child is a plain
resistor or a nested sub-circuit. The base destructor is declared
\texttt{virtual} so deleting a derived object through a base pointer is
well-defined.

\paragraph{Composite pattern.}
Because \texttt{Circuit} itself inherits from \texttt{Component}, a
circuit can contain other circuits. The recursion is implicit:
\texttt{Circuit::resistance} just calls \texttt{resistance} on each child
and relies on virtual dispatch, and \texttt{Circuit::analyze} similarly
calls \texttt{voltageDrop} and \texttt{power} polymorphically.
\texttt{analyze} additionally uses \texttt{dynamic\_cast} to detect when a
child is itself a \texttt{Circuit} so it can recurse with the correct
branch voltage. The same composite property makes \texttt{operator+} and
\texttt{operator*} work --- they wrap clones of both operands inside a
fresh series or parallel circuit.

\paragraph{Smart pointers and RAII.}
The library uses no raw owning pointers. Children are stored as
\texttt{std::unique\_ptr<Component>}, deletion is automatic when the
owning vector goes out of scope, and ownership is unambiguous. Object
construction goes through \texttt{std::make\_unique}. File handles are
managed by \texttt{std::ifstream} / \texttt{std::ofstream}, so streams are
closed automatically even on an exceptional return path.

\paragraph{Virtual cloning.}
Deep copies go through the virtual \texttt{clone()} method on each
subclass, which avoids object slicing when copying a polymorphic
container. \texttt{Circuit::clone} recursively clones every child.

\paragraph{Operator overloading.}
\texttt{Circuit::operator+} and \texttt{Circuit::operator*} let the client
write \texttt{series + parallel} and \texttt{parallel * parallel}, mirroring
the algebraic rules. \texttt{operator<<} is overloaded as a friend free
function on \texttt{Component} and delegates to the virtual \texttt{print}
method, so any node can be streamed with the same syntax.

\paragraph{Templates and the STL.}
The implementation leans on \texttt{std::vector}, \texttt{std::unique\_ptr},
\texttt{std::string}, \texttt{std::ifstream} / \texttt{std::ofstream}, and
\texttt{std::numeric\_limits<double>::infinity()} to represent a
DC-blocking capacitor. \texttt{<cmath>}'s \texttt{std::isinf} and
\texttt{std::isfinite} guard the series sum and parallel reciprocal so an
open branch propagates correctly without producing NaN.

\paragraph{Run-time type identification.}
\texttt{Circuit::analyze} uses \texttt{dynamic\_cast<const Circuit*>} on
each child pointer to decide whether to recurse into a sub-circuit with
its branch voltage; leaf components do not need the recursive call. This
keeps the analysis routine generic over arbitrary nesting depth without
adding a virtual hook to the leaf interface.

\paragraph{Strongly-typed enums.}
\texttt{Circuit::Topology} is declared as \texttt{enum class}, so its
enumerators do not leak into the surrounding scope and cannot be implicitly
converted to integers.

\paragraph{Move semantics.}
Constructors take string parameters by value and \texttt{std::move} them
into the member. \texttt{Circuit::add} accepts a \texttt{std::unique\_ptr}
by value and \texttt{std::move}s it into the vector, making the ownership
transfer explicit at the call site.

\paragraph{Factory pattern.}
\texttt{ComponentFactory::create} maps a type-name string to the matching
concrete object, keeping \texttt{Circuit::load} free of any knowledge of
the leaf classes. Adding a new component type requires touching only the
factory and the new class.

\paragraph{Exceptions.}
\texttt{Circuit::save} and \texttt{Circuit::load} throw
\texttt{std::runtime\_error} when the file cannot be opened, and the
factory throws \texttt{std::invalid\_argument} on an unknown type name,
giving the caller a clear failure path.

\paragraph{Const correctness.}
Methods that do not modify state are marked \texttt{const}. Read-only
access to the child vector is exposed through a \texttt{const\&} getter,
which prevents callers from mutating the internal list while still
allowing efficient iteration.

\section{Summary}

The project meets its functional goals --- building, analyzing,
combining, saving, and loading DC circuits, and reporting per-component
voltage drop and power dissipation --- using a small, idiomatic C++17
codebase. The composite pattern keeps the data model uniform, smart
pointers keep ownership unambiguous, operator overloading gives the user
algebraic syntax, and the factory keeps the loader decoupled from the
concrete element types.

\end{document}
